\documentclass[exercice]{thedocument}%
\usepackage{tabularray}%
\startdatakeys
\source{}
\BO{2026}
\cycle{4}
\competencesDuSocle{CA,MO}
\objectifsApprentissage{C4D1aa,C4D1ab,C4D1ac}
\automatismes{C4D1ab,C4D3af}
\enddatakeys
\begin{document}%
Les deux tableaux suivants traduisent une situation de proportionnalité.
Les compléter en notant les calculs effectués.
\begin{enumerate}
    \item \begin{tblr}{colspec=*3{Q[c,m]},hlines,vlines}
        \<n1;&\<n2;&\<n3;\\
        \<n4;&&
    \end{tblr}
    \item \begin{tblr}{colspec=*3{Q[c,m]},hlines,vlines}
        \<n5;&\<n6;&\<n7;\\
        \<n8;&&
    \end{tblr}
\end{enumerate}
\end{document}
